%!TeX program = xelatex
\documentclass{SYSUReport}
\usepackage{tabularx} % 在导言区添加此行
\usepackage{float}
% 根据个人情况修改
\headl{}
\headc{}
\headr{并行程序设计与算法实验}
\lessonTitle{并行程序设计与算法实验}
\reportTitle{Lab5-基于OpenMP的并行矩阵乘法}
\stuname{林宏宇}
\stuid{23320093}
\inst{计算机学院}
\major{计算机科学与技术}
\date{2026年5月6日}

\begin{document}

% =============================================
% Part 1: 封面
% =============================================
\cover
\thispagestyle{empty} % 首页不显示页码
\clearpage

% % =============================================
% % Part 4： 正文内容
% % =============================================
% % 重置页码，并使用阿拉伯数字
% \pagenumbering{arabic}
% \setcounter{page}{1}

%%可选择这里也放一个标题
%\begin{center}
%    \title{ \Huge \textbf{{标题}}}
%\end{center}

\section{实验目的}
\begin{itemize}
    \item 掌握OpenMP编程的基本流程
    \item 掌握常见的OpenMP编译指令和运行库函数
    \item 分析调度方式对多线程程序的影响
\end{itemize}

\section{实验内容}
\begin{itemize}
 \item 使用OpenMP实现并行通用矩阵乘法
\item 设置不同线程数量（1-16）、矩阵规模（128-2048）、调度方式
\begin{itemize}
    \item 调度方式包括默认调度、静态调度、动态调度
\end{itemize}
\item 根据运行时间，分析程序并行性能
\item 选做：根据运行时间，对比使用OpenMP实现并行矩阵乘法与使用Pthreads实现并行矩阵乘法的性能差异，并讨论分析。
\end{itemize}

\section{实验结果}
\begin{table}[H]
\centering
\caption{默认调度}
\label{表1}
\begin{tabular}{|c|c|c|c|c|c|}
\hline
\multirow{2}{*}{线程数} & \multicolumn{5}{c|}{矩阵规模}                                                                        \\ \cline{2-6} 
 & \multicolumn{1}{l|}{128} & \multicolumn{1}{l|}{256} & \multicolumn{1}{l|}{512} & \multicolumn{1}{l|}{1024} & 2048 \\ \hline
1                    & 4.173ms & 17.872ms & 79.756ms & 621.700ms & 7005.643ms \\ \hline
2                    & 2.271ms & 15.983ms & 57.082ms & 322.121ms & 3729.484ms \\ \hline
4                    & 1.433ms & 8.319ms & 37.781ms & 171.441ms & 1993.081ms \\ \hline
8                    & 1.253ms & 4.607ms & 29.509ms & 94.010ms & 1036.010ms \\ \hline
16                   & 1.828ms & 3.304ms & 16.776ms & 58.566ms & 548.514ms \\ \hline
\end{tabular}
\end{table}

\begin{table}[H]
\centering
\caption{静态调度}
\label{表2}
\begin{tabular}{|c|c|c|c|c|c|}
\hline
\multirow{2}{*}{线程数} & \multicolumn{5}{c|}{矩阵规模}                                                                        \\ \cline{2-6} 
 & \multicolumn{1}{l|}{128} & \multicolumn{1}{l|}{256} & \multicolumn{1}{l|}{512} & \multicolumn{1}{l|}{1024} & 2048 \\ \hline
1                    & 4.165ms & 19.329ms & 79.715ms & 614.283ms & 7192.747ms \\ \hline
2                    & 2.215ms & 16.124ms & 56.845ms & 323.962ms & 4205.540ms \\ \hline
4                    & 1.569ms & 10.893ms & 35.549ms & 171.367ms & 1878.935ms \\ \hline
8                    & 1.153ms & 4.664ms & 26.187ms & 93.836ms & 952.940ms \\ \hline
16                   & 1.554ms & 3.434ms & 16.534ms & 59.603ms & 487.190ms \\ \hline
\end{tabular}
\end{table}

\begin{table}[H]
\centering
\caption{动态调度}
\label{表3}
\begin{tabular}{|c|c|c|c|c|c|}
\hline
\multirow{2}{*}{线程数} & \multicolumn{5}{c|}{矩阵规模}                                                                        \\ \cline{2-6} 
 & \multicolumn{1}{l|}{128} & \multicolumn{1}{l|}{256} & \multicolumn{1}{l|}{512} & \multicolumn{1}{l|}{1024} & 2048 \\ \hline
1                    & 4.585ms & 22.008ms & 87.978ms & 688.606ms & 7356.722ms \\ \hline
2                    & 3.318ms & 17.350ms & 60.035ms & 389.261ms & 3973.898ms \\ \hline
4                    & 2.909ms & 10.383ms & 40.709ms & 206.601ms & 2216.356ms \\ \hline
8                    & 2.022ms & 7.160ms & 28.944ms & 128.815ms & 1111.109ms \\ \hline
16                   & 2.029ms & 4.450ms & 17.734ms & 74.206ms & 583.112ms \\ \hline
\end{tabular}
\end{table}
\section{实验分析}

\subsection{并行性能分析}

由表1--表3的实验数据可以观察到以下并行性能特征：

\textbf{加速比分析}：以默认调度(2048矩阵)为例，单线程运行时间为7005.643ms，2线程为3729.484ms（加速比1.88x），4线程为1993.081ms（加速比3.52x），8线程为1036.010ms（加速比6.76x），16线程为548.514ms（加速比12.77x）。随着线程数增加，程序运行时间显著减少，展现出良好的可扩展性。但由于线程创建和同步开销，以及缓存/内存带宽争用等因素，加速比并非线性增长，尤其在较小矩阵规模下更为明显——例如128矩阵从1线程到16线程的加速比仅为2.28x，远低于理论值。

\textbf{矩阵规模的影响}：矩阵规模越大，并行加速效果越明显。以默认调度16线程为例，128矩阵加速比为2.28x，256矩阵加速比为5.41x，512矩阵加速比为4.75x，1024矩阵加速比为10.62x，2048矩阵加速比为12.77x。总体上，矩阵规模增大使每个线程的计算负载更饱满，有效摊薄了并行开销；在本实验范围内未出现明显的带宽饱和现象。

\subsection{调度方式对比分析}

对比三种调度方式的表现：

\textbf{默认调度与静态调度}：两种调度方式的性能非常接近。OpenMP标准并未强制规定默认调度的具体策略，具体实现可能不同；在GCC等常见实现中，规则的\texttt{for}循环默认采用\texttt{static}调度，因此与显式\texttt{static}在性能上表现接近。本实验中默认调度与静态调度确实接近，这与上述实现习惯一致，也反映出迭代负载较均衡、静态划分开销最低。

\textbf{动态调度}：动态调度在单线程下运行时间略长（2048矩阵：7356.722ms vs 7005.643ms），因为动态调度引入了额外的任务分配开销。但在多线程场景下（如16线程2048矩阵：583.112ms vs 548.514ms），动态调度仍略慢，主要由于调度开销增加；矩阵乘法各迭代负载均匀，因此动态调度的负载均衡优势在本实验中未能体现。在负载不均衡的实际应用场景中，动态调度才能发挥其优势。

\textbf{小结}：对于负载均衡的矩阵乘法任务，推荐使用默认/静态调度以获得最优性能；当各迭代计算量差异较大时，动态调度可通过按需分配任务来避免负载不均。

\subsection{选做：OpenMP与Pthreads性能对比}

为对比OpenMP与Pthreads的性能差异，使用相同的矩阵乘法算法（三层循环，ikj顺序）分别用Pthreads实现并进行相同规模的测试。Pthreads版本手动将矩阵A的行均匀分配给各线程，每个线程独立计算所分配行的C矩阵结果。

\subsubsection{对比数据}

表4展示了Pthreads实现在各线程数和矩阵规模下的运行时间，表5对比了OpenMP（默认调度）与Pthreads在2048规模下的运行时间与加速比。

\begin{table}[H]
\centering
\caption{Pthreads并行矩阵乘法运行时间}
\label{表4}
\begin{tabular}{|c|c|c|c|c|c|}
\hline
\multirow{2}{*}{线程数} & \multicolumn{5}{c|}{矩阵规模} \\ \cline{2-6}
 & 128 & 256 & 512 & 1024 & 2048 \\ \hline
1  & 4.747ms & 24.722ms & 101.874ms & 701.367ms & 8314.071ms \\ \hline
2  & 2.616ms & 16.208ms & 59.925ms & 366.635ms & 4269.847ms \\ \hline
4  & 1.574ms & 11.260ms & 40.019ms & 197.649ms & 2175.121ms \\ \hline
8  & 1.296ms & 4.964ms & 27.837ms & 105.374ms & 1166.436ms \\ \hline
16 & 1.824ms & 5.251ms & 17.381ms & 66.235ms & 617.557ms \\ \hline
\end{tabular}
\end{table}

\begin{table}[H]
\centering
\caption{OpenMP(默认调度)与Pthreads对比(2048规模)}
\label{表5}
\begin{tabular}{|c|c|c|c|c|}
\hline
\multirow{2}{*}{线程数} & \multicolumn{2}{c|}{运行时间(ms)} & \multicolumn{2}{c|}{加速比} \\ \cline{2-5}
 & OpenMP & Pthreads & OpenMP & Pthreads \\ \hline
1  & 7005.643 & 8314.071 & 1.00x & 1.00x \\ \hline
2  & 3729.484 & 4269.847 & 1.88x & 1.95x \\ \hline
4  & 1993.081 & 2175.121 & 3.52x & 3.82x \\ \hline
8  & 1036.010 & 1166.436 & 6.76x & 7.13x \\ \hline
16 & 548.514 & 617.557 & 12.77x & 13.46x \\ \hline
\end{tabular}
\end{table}

\subsubsection{性能分析}

从表5数据可以得出以下结论：

\textbf{绝对性能}：在单线程场景下，OpenMP（7005.643ms）比Pthreads（8314.071ms）快约15.7\%。该差异可能来自编译优化选项、编译参数差异（如向量化/优化级别）、运行时库开销或CPU亲和设置等因素，不能简单归因于编程模型本身。在多线程场景下（如16线程），两者性能趋于接近（548.514ms vs 617.557ms），差距缩小至约11.2\%。

\textbf{加速比与并行效率}：Pthreads的加速比（16线程：13.46x）在数值上高于OpenMP（12.77x），但加速比受单线程基线影响较大，因此更应结合绝对运行时间与并行效率进行判断。从并行效率来看——OpenMP：$E=T_1/(P\times T_P)=7005.643/(16\times548.514)=79.8\%$，Pthreads：$E=T_1/(P\times T_P)=8314.071/(16\times617.557)=84.1\%$——两者均在合理范围内，Pthreads效率略高，这与其手动静态划分在负载均衡场景下调度开销较小一致。

\textbf{编程复杂度}：虽然两者性能接近，但OpenMP在开发效率上具有显著优势。Pthreads需要约100行代码手动管理线程创建、参数传递、任务划分（处理余数行）和线程汇合；而OpenMP仅需一行\texttt{\#pragma omp parallel for}编译制导语句即可实现等价功能。此外，OpenMP的调度策略可通过\texttt{schedule}子句灵活切换（静态、动态等），而Pthreads需要手动修改任务分配逻辑。

\textbf{小结}：对于矩阵乘法这类计算密集型、负载均衡的任务，本实验数据中OpenMP与Pthreads在性能上差异不大（约10\%--16\%），但OpenMP在代码简洁性和可维护性上具有明显优势，是此类场景的首选方案。

\end{document}